\documentclass{ximera}
\title{Smooth manifolds}

\begin{document}

    \begin{abstract} Charts, atlases, smooth structures, and smooth manifolds.   \end{abstract}
    \maketitle

    A reference for this material is Chapter 1 of John M. Lee. Introduction to smooth manifolds. Second edition. Grad. Texts in Math., 218. Springer, New York, 2013. xvi+708pp. ISBN: 978-1-4419-9981-8.

    \begin{definition}[Topological manifold]
        Let $n \in \mathbb{N}$.  An  \emph{$n$-dimensional topological manifold} is a second-countable Hausdorff topological space $M$ such that each point $p \in M$ has an open neighborhood homeomorphic to an open set in $\mathbb{R}^n$.
    \end{definition}
    
    \begin{definition}[Charts and parametrizations]
        Let $M$ be an $n$-dimensional topological manifold, $U \subset M$ an open set, and $\varphi : U \to V$ a homeomorphism with $V \subset \mathbb{R}^n$ open.  The pair $(U, \varphi)$ is called a \emph{chart} and  the pair $(V ,\varphi ^{-1} )$ is called a \emph{parametrization}.  By an abuse of notation, we often call  $\varphi$ a chart.
    \end{definition}
    
    \begin{definition}[Compatible charts]
        Let $M$ be a topological manifold. We say two charts $(U , \varphi )$, $(V, \psi)$  are \emph{compatible} if  the map 
        \[    \psi \circ \varphi ^{-1} : \varphi (U \cap V) \to \psi (U \cap V)      \]
        is a diffeomorphism.
    \end{definition}

    \begin{proposition}[Transfer of compatibility]
        Let $M$ be an $n$-dimensional topological manifold and charts $(U, \varphi )$ and $(V, \psi )$. Assume there is a collection of charts $\{ (U_i, \varphi _i ) \} _{i \in I} $  such that for each $i \in I$, the chart $(U_i, \varphi_i ) $ is compatible with both $(U, \varphi) $ and $(V, \psi)$. If 
    \[    U \cap V \subset \bigcup_{i \in I} U_i ,  \]
    then $U$ and $V$ are compatible. 


    
        \begin{solution}
        Homework.
    \end{solution}
    \end{proposition}



    
    \begin{definition}[Smooth atlas]
        Let $M$ be a  topological manifold. A \emph{smooth atlas} is a collection of compatible charts $\mathcal{A} = \{  (U_i , \varphi_i ) \} _{i \in I } $  with 
        \[       M \subset \bigcup _{i \in I} U_i .                \]
    \end{definition}

    \begin{definition}[Smooth structure]
        Let $M$ be a  topological manifold. A \emph{smooth structure} is a maximal smooth atlas $\mathcal{A}$. In this definition, maximal means that any smooth atlas containing $\mathcal{A}$ equals $\mathcal{A}$. 
    \end{definition}
    
    \begin{proposition}
        Any smooth atlas $\mathcal{A}$ is contained in a unique smooth structure $\mathcal{S}$. Moreover, $\mathcal{S}$ consists precisely of the charts compatible with all charts in $\mathcal{A}$.


    \begin{solution}
        Homework.
    \end{solution}
    \end{proposition}


    \begin{definition}[Smooth manifold]
        A smooth manifold is a pair $(M, \mathcal{S})$ with  $M$ a topological manifold and  $\mathcal{S}$ a smooth structure on $M$.
    \end{definition}

    
    \begin{definition}[Smooth function to $\mathbb{R}^m$]
    Let $(M, \mathcal{S})$ be a smooth manifold and $f : M \to \mathbb{R}^m$ a function. We say $f$ is \emph{smooth} if for any chart $(U , \varphi)$ in $ \mathcal{S}$, the composition
    \[     f \circ \varphi ^{-1} : \varphi (U ) \to \mathbb{R}^m        \]
    is smooth. We denote by $C^{\infty}(M)$ the set of smooth functions $f : M \to \mathbb{R}$. 
    \end{definition}

    \begin{proposition} [Smoothness is local]
    Let $(M, \mathcal{S})$ be a smooth manifold, $\mathcal{A} \subset \mathcal{S}$ a smooth atlas, and  $f : M \to \mathbb{R}^m$ a function. Assume that for any chart $(U, \varphi)$ in $\mathcal{A}$, the composition 
    \[     f \circ \varphi ^{-1} : \varphi (U ) \to \mathbb{R}^m        \]
    is smooth.  Then $f$ is smooth.  

    
    
    \begin{solution}
    Exercise. Use that the change of coordinates is smooth and the chain rule.         
    \end{solution}   
    \end{proposition}


    
    \begin{notation}
    Given a smooth manifold $(M, \mathcal{S})$, by an abuse of noation, whenever we say ``chart'' we mean an element of $\mathcal{S}$.    
    \end{notation}




    \begin{notation}
        While a smooth manifold is technically a pair $(M , \mathcal{S})$, it is often simply denoted by $M$. 
    \end{notation}






\end{document}